\chapter{Evaluationsdesign}

\section{Forschungsdesign}

Zur Beantwortung der Forschungsfrage, wie sich verschiedene \acp{LLM} hinsichtlich ihrer Eignung zur Unterstützung von Stichprobenanalysen im Audit unterscheiden, wird in dieser Arbeit eine vergleichende Evaluation der Modelle durchgeführt. Dazu werden die ausgewählte Modelle auf verschiedene Prüfungsdokumente angewendet. Die Modelle analysieren die Dokumente als auch die Bestätigungen zur Freigabe der Rechnungen, markieren mögliche Auffälligkeiten und begründen ihre Befunde. Die Ergebnisse werden anschließend anhand festgelegter Kriterien bewertet und miteinander verglichen. Damit Unterschiede in den Ergebnissen tatsächlich auf die Modelle zurückzuführen sind, werden alle übrigen Einflussfaktoren konstant gehalten. Alle Modelle erhalten dieselben Rechnungen und dieselben Freigaben sowie denselben Prompt um identische Bedingungen zu gewährleisten.

\section{Auswahl der zu untersuchende Modelle}

Die Auswahl der zu untersuchenden Modelle basiert auf der aktuellen Nutzung von \ac{KI}-Lösungen für die Freudenberg Gruppe. Wie in \ref{sec:ki_at_the_moment} beschrieben nutzt die Freudenberg Gruppe gruppenweit den Microsoft Copilot Premium Version und die dort ausgerollten Modelle von OpenAI und Anthropic. Hierbei werden sowohl die internen Datenschutz- und Compliance Richlinien als auch EU-Richtlinien eingehalten. Ein weiterer Vorteil ist, dass durch die Nutzung der Modelle über den Microsoft Copilot keine weiteren Kosten aufkommen. Zudem ist der Copilot schon im internal Audit ausgerollt, weshalb es keiner weiteren Schulungen und Gewöhnungszeiten für Mitarbeiter braucht.

Auch die Nutzung von Open-Source-Modellen stellt grundsätzlich eine Alternative zu kommerziellen Lösungen dar. Durch den Betrieb auf unternehmensinterner Infrastruktur können Datenschutz- und Compliance-Anforderungen gezielt berücksichtigt werden. Da die Bereitstellung und der Betrieb solcher Modelle jedoch mit erheblichem Infrastruktur- und Implementierungsaufwand verbunden sind, werden Open-Source-Modelle im Rahmen dieser Arbeit nicht betrachtet.

Der Microsoft Copilot ermöglicht es den Auditieren auf die Modelle Claude Sonnet, Claude Opus, GPT 5.6 Sol (Schnelle Analyse), GPT 5.6 Sol (Tiefere Analyse) zuzugreifen. Genauere Details, bei welcher Anfrage der Copilot auf GPT 5.6 Sol, GPT 5.6 Terra oder GPT 5.6 Luna zugreift sind nicht bekannt. \todo{nochmal schauen welceh modelle es jetzt gibt}. Die vier Modelle gehören laut dem Artifical Analysis Intelligence Index zu den aktuell mit am besten gewerteten kommerziellen \ac{KI}-Modellen und sind somit für den Anwendungsfall bei der \ac{FCo} geeignet (vgl. \cite{ArtificialAnalysis2026}).

\section{Entwicklung der Testdatensätze}

Für die Evaluation der betrachteten \ac{KI}-Modelle ist eine Datenbasis erforderlich, die den realen Prüfungskontext möglichst genau abbildet. Die Verwendung originaler Rechnungsbelege aus dem Internal Audit der Freudenberg-Gruppe ist ausgeschlossen, da diese personenbezogene sowie geschäftskritische Informationen enthalten und daher im Rahmen dieser Evaluierung nicht herangezogen werden dürfen. Um dennoch eine hinreichend umfangreiche und belastbare Testbasis zu schaffen, werden die Testdatensätze mithilfe generativer \ac{KI} erzeugt. Die künstlich erstellten Rechnungen orientieren sich in Struktur, Umfang und Informationsgehalt an realen Rechnungen und enthalten dieselben prüfungsrelevanten Kriterien wie die im Audit untersuchten Dokumente. Nach der Generierung der Rechnungen werden gezielt Fehler in die Dokumente eingebracht. Dieser Ansatz orientiert sich an Benchmark-Studien zur Evaluation von \acp{LLM} im Finanzkontext, bei denen vorhandene Finanzdokumente systematisch durch künstlich eingefügte Fehler verändert werden, um die Fähigkeit der Modelle zur Identifikation von Inkonsistenzen zu untersuchen (vgl. \cite[S. 1]{Panda2026FinVerBench}). Neben den Rechnungen werden auch Screenshots von Mails zur Freigabe der einzelnen Rechnungen erstellt. Auch hier werden gezielt einzelen Fehler in die Datensätze eingebaut. Der Testdatensatz umfasst somit sowohl fehlerhafte als auch vollständig korrekte Rechnungen sowie die zugehörigen Freigaben. Für die vorliegende Arbeit wird eine Fehlerquote von 24\,\% festgelegt, was bei einem Datensatz von 100 Rechnungen insgesamt 24 fehlerhaften Belegen entspricht. Nach internen Angaben ist die Quote für fehlerhafte oder unvollstänidge Rechnungen bei 18\% (vgl. internes Meeting David Goth). Durch eine Fehlerquote von 24 \% ist die Anzahl fehlerhafter Rechnungen ausreichend groß, um die Erkennungsleistung der Modelle repräsentativ zu untersuchen \todo{Quelle? ansonsten auch David}. Gleichzeitig bildet die Mehrheit an korrekten Rechnungen den realen Prüfungskontext realitätsnah ab. Bei der Fehlererzeugung wird bewusst auf eine breite Streuung unterschiedlicher Fehlertypen geachtet. Die Fehler werden in sechs Gruppen unterschieden. Die größte Gruppe bilden 7 Fehler, die typische Betrugsmuster abbilden. Sie zeigen den eigentlichen Nutzen eines solchen Systems im Audit. Bei 5 Fehlern liegt das Problem im Freigabeprozess, wenn eine Rechnung ohne die erforderliche Freigabe oder durch eine nicht berechtigte Person genehmigt wird. Weitere 4 Fehler betreffen den Abgleich mit den Stammdaten, bei denen die Angaben auf der Rechnung von den hinterlegten Lieferantendaten abweichen, zum Beispiel bei Anschrift oder Bankverbindung. Bei 4 Fehlern handelt es sich um einfache Rechen-, bzw. Zahlenfehler, die als Basisfälle gelten. Ebenfalls 3 Fehler betreffen ein falsches Rechnungsdatum oder die Zuordnung zur falschen Periode, sodass der Aufwand im falschen Zeitraum gebucht wird. Die letzten 2 Fehler betreffen eine fehlende Rechnungsnummer. Tabelle~\ref{tab:fehlergruppen} fasst die Fehlergruppen zusammen.

\todo{genaue zahlen begründen: 100 rechnunge, 24 \%, 24 fehler, fehler 7, 54, 3, 3, 2}

\begin{table}[htbp]
  \centering
  \caption{Verteilung der eingebauten Fehler auf die Fehlergruppen}
  \label{tab:fehlergruppen}
  \begin{tabular}{cl}
    \hline
    \textbf{Anzahl} & \textbf{Fehlergruppe}\\
    \hline
    7 & Betrugsmuster \\
    5 & Fehler in der Rechnungsfreigabe\\
    4 & Abgleich mit Stammdaten \\
    4 & Rechnerische Fehler und Basisfälle \\
    3 & Datum und Periode \\
    2 & Fehlende Rechnungsnummer \\
    \hline
    24 & \textbf{Gesamt} \\
    \hline
  \end{tabular}
\end{table}

\section{Entwicklung eines geeigneten Prompts}

Die Qualität der von einem \ac{LLM} erzeugten Ausgabe wird maßgeblich durch die Gestaltung des Prompts beeinflusst (vgl. \cite{ICLR2024_6c0e99d7}). Prompt Engineering bezeichnet die systematische Gestaltung, Strukturierung und Optimierung von Eingabeanweisungen (Prompts) für \acp{LLM}. Ziel ist es, durch die gezielte Auswahl von Aufgabenbeschreibungen, Kontextinformationen, Beispielen und Formatierungen das gewünschte Modellverhalten zu erzeugen und die Qualität der generierten Ergebnisse zu verbessern (vgl. \cite[S. 1]{schulhoff2025promptreportsystematicsurvey}). Daher wird der Prompt auf Basis etablierter Prompt-Engineering-Techniken entwickelt. Um die Vergleichbarkeit der untersuchten Modelle sicherzustellen, wird auf eine modellspezifische Optimierung der Prompts verzichtet. Stattdessen werden für alle Modelle identische Prompts verwendet. Der Prompt enthält klare Anweisungen zur Analyse der Rechnungen, zur Identifikation potenzieller Auffälligkeiten sowie zur nachvollziehbaren Begründung der Ergebnisse.

Die inhaltlichen Anforderungen an den Prompt ergeben sich aus dem bestehenden Vorgehen der Stichprobenanalyse bei Freudenberg (vgl. Abschnitt~\ref{sec:stichprobenanalyse_freudenberg}). Der Prompt bildet die Prüfschritte ab, die ein Prüfer bei einer Eingangsrechnung durchführt. Zuerst die Extraktion der Rechnungsdaten, dann den Abgleich mit der Referenztabelle, sowie die rechnerische und formale Prüfung und die Suche nach Betrugsindikatoren. Abschließend noch die Prüfung des Freigabeprozesses und der Berechtigung des Genehmigers. Das Ergebnis soll als Arbeitspapier für die anschließende menschliche Prüfung dienen. Das Modell unterstützt den Prüfer also, trifft aber lediglich eine abschließende Empfehlung für weiteres Vorgehen und übernimmt nicht den vollstädnigen Prüfungsprozess. \todo{Freudenberg klar machen -> Freudenberg CO. KG}

Der Prompt besteht aus zwei Teilen. Die Instruktionen sind dauerhaft im Agenten hinterlegt und legen Rolle, Ziel, Quelen, Regeln, Prüfschritte und Ausgabeformat fest. Der Kurzprompt wird für jeden Prüffall mit gleichem Wortlaut in den Chat eingegeben. Dieser Prompt führt den Agent aus und benennt die vorliegenden Unterlagen.. Tabelle~\ref{tab:promptaufbau} zeigt die Abschnitte der Instruktionen. Der vollständige Prompt unterteilt in Agenteninstructions und Kurzprompt befindet sich in Anhang % ref anhang}.

Bei der Gestaltung des Prompts wurden mehrere etablierte Techniken angewendet.
Zu Beginn erhält das Modell eine klar definiertes Ziel und eindeutige Rolle als Prüfagent des internen Audits. Die Rolle gibt dem Modell den fachlichen Rahmen der Aufgabe vor (vgl. \cite{microsoft2026declarative}).

Für jede Eingabe ist festgelegt, welche Funktion sie hat. Die Rechnung liefert die Istwerte, die Referenztabelle die Sollwerte und die Freigabematrix die Berechtigungen. Die Freigabeunterlagen gelten ausschließlich als Evidenz und sind nicht automatisch gültig. Durch diese Quellenhierarchie entscheidet das Modell bei widersprüchlichen Angaben nicht selbst, welcher Quelle es folgt. Die einzelnen Wissensquellen werden in den Instructions ausdrücklich definiert und erklärt (vgl. \cite{microsoft2026declarative}).

Die Aufgabe ist in elf Prüfschritte mit fester Reihenfolge zerlegt. Die Aufteilung in einzelne, klar getrennte Einheiten verringert Mehrdeutigkeiten und verhindert, dass das Modell Aufgaben zusammenführt oder neu interpretiert (vgl. \cite{microsoft2026declarative}). Zusätzlich wird die Prüfung dadurch nachvollziehbar, da jeder Schritt im Ergebnis sichtbar ist.

\acp{LLM} neigen dazu, fehlende Informationen durch plausibel klingende, aber falsche Inhalte zu ersetzen \cite[S. 1]{huang2025survey}. Der Prompt verbietet deshalb Annahmen und Ergänzungen und gibt feste Statuswerte wie „nicht vorhanden“, „nicht lesbar“ und „nicht prüfbar“ vor. Verdachtsmomente dürfen nur als „Hinweis auf“ formuliert werden, da die Feststellung eines Betrugs dem Prüfer vorbehalten bleibt (vgl. \cite{anthropic2026hallucinations}).

Die Ausgabe folgt einer festen Struktur mit vorgegebenen Feldern und Statuswerten. Ohne vorgegebenes Ausgabeformat kann es zu inkonsistentem Verhalten zwischen verschiedenen Modellen kommen (vgl. \cite{microsoft2026declarative}). Das feste Format erleichtert die menschliche Prüfung und macht die Ergebnisse der Modelle direkt vergleichbar. Die Gesamtbewertung der Freigabe, die Auffälligkeiten mit Risikostufe sowie die Angabe zur menschlichen Nachprüfung bilden die Grundlage für die Auswertung.

Im Prüfkontext ist ein übersehener Fehler kritischer als ein zusätzlicher Hinweis, da eine fehlerhafte Rechnung sonst unentdeckt bleibt. Der Prompt gibt diese Gewichtung ausdrücklich vor. \todo{wirkt iwie random}


Der Prompt wurde in mehreren Versionen entwickelt. Getestet wurde dabei ausschließlich mit einem separaten , der nicht in die spätere Evaluation eingeht. Würde der Prompt anhand der Testfälle verbessert, wären die gemessenen Ergebnisse zu positiv \cite{Q8}.

Die erste Version enthielt ausführliche Definitionen, Beispiele und Statusbeschreibungen. % TODO: Zeichenanzahl und Limit der Plattform eintragen
Mit etwa [X] Zeichen überschritt sie das Zeichenlimit der eingesetzten Agentenplattform von [Y] Zeichen. Zudem berücksichtigen \acp{LLM} Informationen in langen Kontexten weniger zuverlässig \cite{Q9}. Ziel der Überarbeitung war daher, dem Modell nur die für die Aufgabe notwendigen Informationen bereitzustellen \cite{Q1}. Die zweite Version fasst die Regeln stichpunktartig zusammen, ohne Prüfschritte zu streichen. In der dritten Version wurden Widersprüche zwischen Instruktionen und Kurzprompt beseitigt. Tabelle~\ref{tab:promptversionen} fasst die Änderungen zusammen.

\begin{table}[htbp]
\centering
\begin{tabular}{|p{1.2cm}|p{6.5cm}|p{6cm}|}
\hline
\textbf{Version} & \textbf{Änderung} & \textbf{Grund} \\
\hline
v1 & Ausführliche Fassung mit Platzhaltern für die Eingaben, Beispielen und Statusdefinitionen & Ausgangsversion \\
\hline
v2 & Kompakte Fassung, Regeln stichpunktartig zusammengefasst & Zeichenlimit der Plattform, Fokus auf notwendige Informationen \\
\hline
v3 & Dateinamen und Ablageorte der Quellen ergänzt, Freigabematrix im Kurzprompt ergänzt, Rundungstoleranz wiederhergestellt, Begriffe vereinheitlicht, Prüfung mehrstufiger Freigaben entfernt & Widersprüche zwischen Instruktionen und Kurzprompt, Verlust von Regeln durch die Kürzung, mehrstufige Freigaben im Prüfprozess nicht vorgesehen \\
\hline
\end{tabular}
\caption{Versionen des Prompts}
\label{tab:promptversionen}
\end{table}

Nach Abschluss der Entwicklung wurde der Prompt eingefroren und für alle Modelle unverändert verwendet. Die Rahmenbedingungen der Testdurchführung, etwa ein neuer Chat je Rechnung, werden in Abschnitt~\ref{sec:durchfuehrung} beschrieben.

\section{Nutzwertanalyse}
Die Eignung der untersuchten \ac{KI}-Modelle für den definierten Audit-Use-Case kann nicht anhand eines einzelnen Bewertungskriteriums bestimmt werden. Vielmehr müssen sowohl quantitative, als auch qualitative Kriterien berücksichtigt werden um ein umfangreiches Ergebnis zu erlangen. Die Bewertung der Modelle erfordert somit die gleichzeitige Betrachtung mehrerer, teilweise konkurrierender Kriterien und stellt ein mehrdimensionales Entscheidungsproblem dar (vgl. \cite[S. 116 f.]{kozlowska2022}). Zur strukturierten und nachvollziehbaren Berücksichtigung dieser unterschiedlichen Kriterien wird in dieser Arbeit eine Nutzwertanalyse durchgeführt. Die Nutzwertanalyse ermöglicht es, verschiedene Bewertungskriterien zu gewichten und zu einem Gesamtnutzenwert zu aggregieren. Dadurch können die untersuchten \ac{KI}-Modelle objektiv und transparent hinsichtlich ihrer Eignung für den Audit-Use-Case verglichen werden (vgl. \cite[S. 79 f.]{kuehnapfel2021}).


\section{Entwicklung einer Bewertungsmatrix}
Um die Evaluation basierend auf einer Nutzwertanalyse durchzuführen, muss eine Bewertungsmatrix mit mehreren quantitativen und qualitativen Kriterien entwickelt werden (vgl. \cite[S.3]{liang2023holisticevaluationlanguagemodels}).
Bei der Entwicklung der Bewertungsmatrix werden die Anforderungen für den Anwendungsfall der Corporate Audit Abteilung der \ac{FCo} berücksichtigt. Die Bewertungsmatrix umfasst drei Hauptkriterien, die für die Eignung der Modelle für den definierten Use Case von zentraler Bedeutung sind. Diese Kriterien sind die Erkennungsleistung der Modelle, die Qualität der Begründungen und die Konsistenz der Ergebnisse. Sind alle drei Kriterien hinreichend erfüllt, können die Modelle als geeignet für den Einsatz im Audit-Use-Case betrachtet werden.

\todo{wissenschaftlich warum diese Kriterien}


Als zentrales quantiatives Kriterium wird die Erkennugsleistung der Modelle herangezogen. Hierbei werden die Kennzahlen Precision und Recall berechnet. Der Precision-Wert gib an, wie viele der vom Modell als fehlerhaft identifizierten Rechnungen tatsächlich fehlerhaft sind. Er berechnet sich aus dem Quotient aller richtig als fehlerhafte klassifizierten Rechnungen, den sogenannten True Positivies, un der Summe aller als fehlerhaft klassifizierten Rechnungen, also der Summe aus True Positives und False Positives. Der Recall-Wert gibt an, wie viele der tatsächlich fehlerhaften Rechnungen vom Modell korrekt identifiziert wurden. Er berechnet sich aus dem Quotient der True Positives, und der Summe aller tatsächlich fehlerhaften Rechnungen, also der Summe aus True Positives und False Negatives.

Als weiteres Kriterium wird die Qualität der Begründungen der Modelle bewertet. Die qualitative Bewertung der Begründung erfolgt durch 2 Senior Auditoren der \ac{FCo}. Die Auditoren bewerten die Begründungen der Modelle anhand von den drei Kriterien Korrektheit, Nachvollziehbarkeit und Halluzination. Jedes Kriterium wird auf einer Skala von 1 bis 5 bewertet, wobei 1 für "sehr schlecht" und 5 für "sehr gut" steht. Anschließend wird der Mittelwert der Bewertung der Auditoren gebildet und als Nutzwert der Qualität der Begründung herangezogen.
Als letztes Kriterium wird die Konsistenz der Modelle bewertet. Dazu wird jede Rechnung mit identischem Prompt fünfmal von jedem Modell analysiert. Anschließend werden alle unterschiedlichen Befunde erfasst, die über die fünf Durchläufe hinweg gemeldet wurden. Für jeden dieser Befunde wird ermittelt, in wie vielen Durchläufen er auftaucht. Ein Befund, der in allen fünf Durchläufen gemeldet wird, erhält somit den Wert 1, ein Befund, der nur in einem Durchlauf auftaucht, den Wert 0,2. Die Konsistenz ergibt sich aus dem Mittelwert dieser Anteile über alle Befunde und wird in Prozent angegeben. Liefert ein Modell in allen Durchläufen exakt dieselben Befunde, erreicht es eine Konsistenz von 100\,\%. Je häufiger sich die gemeldeten Befunde zwischen den Durchläufen unterscheiden, desto geringer fällt der Wert aus.

\section{Gewichtung der Bewertungsmatrix}
\todo{Gewichtung klar, aber man kann auch anders gewichten - siehe 1. PA}

Um die drei Hauptkriterien zu einem Gesamtergebnis zusammenzuführen, werden diese entsprechend ihrer Bedeutung für den Audit-Use-Case gewichtet. Die Gewichtung erfolgt auf Basis der Anforderungen der Corporate Audit Abteilung der \ac{FCo} und wird in Tabelle~\ref{tab:gewichtung} dargestellt.

Die Erkennungsleistung erhält mit 50\,\% das höchste Gewicht, da die zuverlässige Identifikation fehlerhafter Rechnungen die Kernaufgabe der Stichprobenanalyse darstellt. Innerhalb dieses Kriteriums wird der Recall mit 30\,\% höher gewichtet als die Precision mit 20\,\%. Ein False Positive also ein Fehlalarm, wird vom Auditor im Rahmen der weiteren Prüfung erkannt und verursacht lediglich zusätzlichen Prüfaufwand. Wird ein Fehler hingegen vom Modell nicht gemeldet, False Negative, besteht die Gefahr, dass auch der Auditor diesen übersieht, da er sich auf das Ergebnis des Systems verlässt.

Die Qualität der Begründungen wird mit 30\,\% gewichtet. Auditoren müssen nachvollziehen können, wie ein Ergebnis zustande gekommen ist, um die Befunde prüfen und dokumentieren zu können. Die Korrektheit erhält hierbei mit 12\,\% das höchste Gewicht, da eine fachlich falsche Argumentation auch bei richtig erkanntem Fehler zu einer fehlerhaften Dokumentation führen kann. Die Nachvollziehbarkeit wird mit 10\,\% und die halluzination mit 8\,\% gewichtet. Erfundene Angaben sind im Audit zwar besonders kritisch, werden aber teilweise bereits über die Korrektheit erfasst.

Die Konsistenz erhält mit 20\,\% das geringste Gewicht. Reproduzierbare Ergebnisse sind zwar eine Voraussetzung dafür, dass Auditoren einem Modell vertrauen können. Eine hohe Konsistenz sagt allein jedoch nichts über die Richtigkeit der Ergebnisse aus, da auch ein Modell, das dauerhaft dieselben Fehler übersieht, konsistent arbeitet. Tabelle~\ref{tab:gewichtung} fasst die Gewichtung zusammen.

\begin{table}[h]
    \centering
    \begin{tabular}{llll}
        \hline
        \textbf{Hauptkriterium} & \textbf{Gewicht} & \textbf{Teilkriterium} & \textbf{Gewicht} \\
        \hline
        Erkennungsleistung & 50\,\% & Recall & 30\,\% \\
                           &        & Precision & 20\,\% \\
        Qualität der Begründungen & 30\,\% & Korrektheit & 12\,\% \\
                           &        & Nachvollziehbarkeit & 10\,\% \\
                           &        & Faktentreue & 8\,\% \\
        Konsistenz         & 20\,\% & Konsistenz & 20\,\% \\
        \hline
    \end{tabular}
    \caption{Gewichtung der Bewertungsmatrix}
    \label{tab:gewichtung}
\end{table}

Da die Kriterien auf unterschiedlichen Skalen erfasst werden, müssen sie vor der Verrechnung vereinheitlicht werden. Precision, Recall und Konsistenz liegen bereits als Prozentwerte vor. Für die Teilkriterien der Begründungsqualität wird zunächst der Mittelwert über alle Auditoren gebildet. Dieser wird anschließend über $e = \frac{x-1}{4} \cdot 100$ in einen Wert zwischen 0 und 100 umgerechnet, sodass die schlechteste Bewertung den Wert 0 und die beste Bewertung den Wert 100 erhält. Somit sind alle Kriterien auf einer gemeinsamen Skala. Der Nutzwert eines Modells ergibt sich dann als Summe der gewichteten Einzelwerte:

\begin{equation}
    N_j = \sum_{i=1}^{6} w_i \cdot e_{ij}
\end{equation}

Dabei steht $N_j$ für den Nutzwert des Modells $j$, $w_i$ für das Gewicht des Teilkriteriums $i$ und $e_{ij}$ für den normierten Wert des Modells $j$ im Teilkriterium $i$.

Da jede Gewichtung eine gewisse Subjektivität enthält, wird im Rahmen der Auswertung zusätzlich eine Sensitivitätsanalyse durchgeführt. Hierbei werden die Gewichte der Hauptkriterien verändert, um zu prüfen, ob sich die Rangfolge der Modelle dadurch verschiebt. Bleibt die Rangfolge stabil, kann das Ergebnis als robust betrachtet werden.

\todo{Sensitivitätsanalyse? Brauch ich die?}